\begin{theorem}[The six-spin flux operation in the actual torus contraction]%
    \label{thm:peps_torus_two_plaquette_global_flux_move}
    \lean{TNLean.PEPS.IsGIsometric.exists_unitary_torusTwoPlaquetteGlobalFluxMove}
    \leanok
    \uses{thm:peps_two_plaquette_torus_geometry,
        thm:peps_torus_two_plaquette_flux_holonomy,
        thm:peps_regular_two_cycle_global_flux_move}
    Let the horizontal and vertical periods be at least four and three,
    respectively, and let the original four-leg tensor be $G$-isometric for
    the regular virtual action. In the six-site region
    $R=\{0,1,2\}\times\{0,1\}$, let $u_g$ insert $g$ on the upward bond
    at horizontal coordinate zero, and let $u'_g$ insert $g$ on both upward
    bonds at horizontal coordinates zero and one. All other internal bonds
    carry the identity.
    There exists one unitary $W$ on the six original spins whose extension
    $\widehat W=W\otimes I$ is a global unitary, such that
    \begin{align}
        \widehat W\,\Psi(u_g;u)=\Psi(u'_g;u)
    \end{align}
    for every $g\in G$ and every common assignment $u$ on exterior and crossing
    bonds. Both vectors are actual globally contracted torus states, with
    the stated internal insertions and the same remaining bond operators.
    The unitary is chosen before $g$ and $u$.
    This is the regular-representation six-site specialization of the local
    flux-moving construction in \cite[Theorem~6.16]{Schuch2010PEPS}; no
    parent-Hamiltonian ground-space membership is asserted.
\end{theorem}

\begin{proof}\leanok
    Express the four-leg tensor in incident-edge coordinates; its
    $G$-isometry is preserved. The fixed spanning path of the actual rectangle
    leaves exactly the leftmost and middle upward bonds as the chosen cycle
    bonds. Apply the global physical two-cycle operation with this fixed tree,
    root, and pair of bonds. The literal insertion assignments equal its
    tree-cycle assignments. Substitution yields the asserted global state
    identity for every common exterior and crossing assignment. Extension by
    the identity preserves unitarity.
\end{proof}
